
\section{Kỹ Thuật Sử Dụng Và Nguồn Gốc}

\begin{quote}\itshape Mục đích: phân biệt rõ kỹ thuật \textbf{kế thừa từ bài báo} vs \textbf{đề xuất mới của đề tài} để tránh đạo văn.\end{quote}

\vspace{0.5em}\hrule\vspace{0.5em}


\subsection{Kỹ Thuật Kế Thừa (Có Cite Bài Báo Gốc)}

\begin{table}[htbp]
\centering\renewcommand{\arraystretch}{1.2}
\resizebox{\textwidth}{!}{%
\begin{tabular}{|l|l|l|}
\hline
\textbf{Kỹ thuật} & \textbf{Kế thừa từ} & \textbf{Cách sử dụng trong đề tài} \\
\hline
\textbf{Kiến trúc U-Net} & Ronneberger 2015 [1] & Làm nền tảng encoder-decoder; không thay đổi cấu trúc cơ bản \\
\hline
\textbf{Attention Gate} & Oktay 2018 [2] & Kế thừa công thức tính hệ số chú ý $\alpha\^{}l$; mở rộng thành Conditioned Attention bằng cách thêm điều kiện hóa câu nhắc \\
\hline
\textbf{Kiến trúc SAM-Med2D} & Cheng 2023 [3] & Dùng làm mô hình đối sánh (benchmark); không tích hợp vào hệ thống đề xuất \\
\hline
\textbf{Backbone MobileNetV4} & Qin 2024 [4] & Dùng trọng số pretrained ImageNet-1K; fine-tune 2 giai đoạn trên BTXRD \\
\hline
\textbf{Thuật toán GradCAM} & Selvaraju 2019 [5] & Kế thừa công thức tính $\alpha\_k\^{}c$ và $L\_\text{GradCAM}$; ứng dụng vào cơ chế cứu hộ câu nhắc (ứng dụng mới) \\
\hline
\textbf{Dataset BTXRD} & Nguyen 2024 [6] & Toàn bộ dữ liệu huấn luyện và đánh giá \\
\hline
\textbf{YOLOv11s} & Ultralytics 2024 & Dùng pretrained COCO; fine-tune trên annotations nhiễu R/L của BTXRD \\
\hline
\textbf{Roboflow} & Roboflow Inc. & Nền tảng gán nhãn bounding box; không có đóng góp kỹ thuật \\
\hline
\end{tabular}%
}
\end{table}

\vspace{0.5em}\hrule\vspace{0.5em}


\subsection{Kỹ Thuật Đề Xuất Mới (Đóng Góp Gốc Của Đề Tài)}

\begin{table}[htbp]
\centering\renewcommand{\arraystretch}{1.2}
\resizebox{\textwidth}{!}{%
\begin{tabular}{|l|l|l|}
\hline
\textbf{Kỹ thuật} & \textbf{Mô tả} & \textbf{Cơ sở lý luận} \\
\hline
\textbf{Prompt Spatial Gate (PSG)} & Tích hợp bản đồ nhiệt vào bộ mã hóa qua phép nhân tăng cường chọn lọc & Chưa có trong bài báo nào; lấy cảm hứng từ cơ chế gating nhưng thiết kế cho U-Net encoder \\
\hline
\textbf{Conditioned Attention Decoder (CAD)} & Điều kiện hóa tín hiệu gating bằng câu nhắc với trọng số tin cậy giảm dần theo tầng & Mở rộng Attention Gate [2]; thêm cơ chế fusion câu nhắc qua confidence score và trọng số w ∈ {1.0, 0.7, 0.4, 0.2} \\
\hline
\textbf{Plateau Heatmap} & Gán 1.0 đồng đều bên trong box + Gaussian blur viền k=31 & Lấy cảm hứng từ Gaussian 2D [lý thuyết] và bản đồ nhiệt GradCAM [5]; biến thể thực tiễn phù hợp U-Net \\
\hline
\textbf{IPR (Iterative Prompt Refinement)} & Vòng lặp tinh chỉnh câu nhắc dùng tâm hình học mặt nạ dự đoán, tối đa 3 vòng & Hoàn toàn mới; không kế thừa từ bài báo nào \\
\hline
\textbf{Pipeline GradCAM Rescue} & Bộ kiểm duyệt 3 tiêu chí $\rightarrow$ GradCAM trích điểm neo $\rightarrow$ IPR cứu hộ & Ứng dụng mới của GradCAM [5] trong ngữ cảnh sửa câu nhắc sai; kết hợp với IPR là đóng góp gốc \\
\hline
\textbf{Two-stage fine-tuning MobileNetV4} & Giai đoạn 1: đóng băng backbone 5 epoch; Giai đoạn 2: mở khóa toàn bộ & Chiến lược fine-tuning chuẩn; áp dụng cụ thể cho bài toán phân lớp BTXRD \\
\hline
\end{tabular}%
}
\end{table}

\vspace{0.5em}\hrule\vspace{0.5em}


\subsection{Phân Biệt Kế Thừa Vs Đề Xuất Mới}

\begin{verbatim}
Kế thừa nguyên si:
  U-Net architecture ────────────────── [1]
  Attention Gate formula ─────────────── [2]
  GradCAM formula (α_k, L_GradCAM) ──── [5]

Kế thừa + mở rộng:
  Attention Gate → Conditioned Attention (thêm câu nhắc) ── dựa trên [2]
  Gaussian heatmap → Plateau Heatmap (biến thể thực tiễn) ── cảm hứng từ lý thuyết

Đề xuất gốc (không kế thừa):
  Prompt Spatial Gate (PSG)
  IPR (Iterative Prompt Refinement)
  Pipeline GradCAM Rescue (ứng dụng GradCAM vào sửa câu nhắc sai)
\end{verbatim}
